% ---------------------------------------------------------------------------
% Fachmodul Endliche Automaten (Kristian Popov)
% ---------------------------------------------------------------------------

\addsectionframe{Fachmodul: Endliche Automaten}

% ---------------------------------------------------------------------------
\begin{itemframe}{Endliche Automaten -- Überblick}
  \item \textbf{DFA} und \textbf{NFA} selbst definieren
  \item \textbf{Rechtslineare Grammatik} $\rightarrow$ \textbf{Automat}
  \item Zwei Sichten: \textbf{Übergangstabelle} und \textbf{Zustandsdiagramm}
  \item \textbf{Stapeltest} für mehrere Wörter gleichzeitig
  \item \textbf{Schritt-Historie}: vorwärts, rückwärts, zu einem früheren Schritt
  \item \textbf{Übungen} mit Rückmeldung: korrekt, falsch, Gegenbeispiel

  \note{
          \textbf{Person: Kristian}

          \textbf{Zeit: 1,5 min}

          \begin{itemize}
            \item Automat selbst definieren: Typ, Alphabet, Zustände, Startzustand, Endzustände,
            Übergänge und testen ihn anschließend mit einem Testwort.
            \item \textbf{Grammatik} $\rightarrow$ \textbf{Automat}: Man kann eine RL Grammatik eingeben und\ldots.
            Jedes Nichtterminal wird Zustand, jede Produktion wird Übergang
            \item \textbf{Zwei Sichten}: Tabelle und Diagramm zeigen denselben Schritt, immer synchron
            \item \textbf{Stapeltest}: Prüft mehrere Wörter auf einmal.
            \item \textbf{Schritt-Historie}: Vorwärts rückwärts oder zu einem bestimmten Schritt springen.
            \item \textbf{Übungen}: Übungen werden automatisch korrigiert.
            Dabei prüft das System das Verhalten des Automaten und bei Fehlern zeigt direkt das erste falsche Wort als Gegenbeispiel an.
          \end{itemize}
  }
\end{itemframe}

% ---------------------------------------------------------------------------
\begin{fullframe}{Simulationsschritt}
  \begin{center}
  \resizebox{0.98\textwidth}{!}{%
    \begin{tikzpicture}[
        >=Stealth,
        stepbox/.style={draw=thwsblue, very thick, rounded corners=3pt, fill=thwsblue!5,
                        minimum width=32mm, minimum height=10mm, align=center, font=\small},
        exitbox/.style={draw=thwsblue!55, thick, rounded corners=3pt, fill=white,
                        minimum width=36mm, minimum height=9mm, align=center,
                        font=\scriptsize, text=thwsblue},
        histbox/.style={draw=thwsorange, very thick, rounded corners=4pt, fill=thwsorange!8,
                        minimum width=62mm, minimum height=12mm, align=center, font=\small},
        viewbox/.style={draw=thwspetrol, very thick, rounded corners=3pt, fill=thwspetrol!8,
                        minimum width=40mm, minimum height=10mm, align=center, font=\small}
      ]

      % --- (1) der Kern: immer dieselbe Schleife
      \node[stepbox] (in)   at (-4.9, 2.4) {aktive Zustände\\+ nächstes Symbol};
      \node[stepbox] (look) at ( 0.0, 2.4) {Übergänge\\nachschlagen};
      \node[stepbox] (out)  at ( 4.9, 2.4) {neue aktive\\Zustände};

      \draw[->, very thick, thwsblue] (in)   -- (look);
      \draw[->, very thick, thwsblue] (look) -- (out);

      \draw[->, thick, thwsblue!70] (out.north) -- ++(0, 0.6)
        -- node[font=\scriptsize, above, text=thwsblue] {nächstes Zeichen} ($(in.north)+(0, 0.6)$)
        -- (in.north);

      % --- (2) wie ein Lauf endet: eigene Kaesten mit Pfeil statt freier Text
      \node[exitbox] (dead) at (0.0, 0.85) {kein Übergang:\\Wort abgelehnt};
      \node[exitbox] (fin)  at (4.9, 0.85) {Wort zu Ende:\\Endzustand aktiv?};

      \draw[->, thick, thwsblue!55] (look.south) -- (dead.north);
      \draw[->, thick, thwsblue!55] (out.south)  -- (fin.north);

      % --- (3) der Schritt wird angehaengt, nicht ueberschrieben
      \node[histbox] (hist) at (0, -1.05)
        {Schritt wird an die \textbf{Historie} angehängt\\[0.4ex]
         {\scriptsize der vorherige Schritt bleibt unverändert}};

      \draw[->, very thick, thwsorange] (out.east) -- ++(1.25, 0) |- (hist.east);

      % --- (4) eine Quelle, zwei Sichten
      \node[viewbox] (tab)  at (-3.0, -2.95) {Übergangstabelle\\{\scriptsize markierte Zelle}};
      \node[viewbox] (diag) at ( 3.0, -2.95) {Zustandsdiagramm\\{\scriptsize markierte Kante}};

      \draw[->, very thick, thwspetrol] (hist.south) -- (tab.north);
      \draw[->, very thick, thwspetrol] (hist.south) -- (diag.north);

      \node[font=\scriptsize, text=thwspetrol, fill=white, inner sep=0.6mm] at (0, -2.1)
        {beide lesen denselben Schritt};
    \end{tikzpicture}%
  }
  \end{center}

  \note{
        \textbf{Person: Kristian}

        \textbf{Zeit: 1 min}

        \begin{itemize}
          \item \textbf{Aktive Zustände} + \textbf{nächstes Symbol}
          $\rightarrow$ passende Übergänge suchen.
          \item Daraus entstehen die neuen aktiven Zustände
          \item Nächstes Zeichen beginnt der Ablauf wieder von vorne.
          \item Zwei Ausgänge: kein Übergang $\rightarrow$ abgelehnt.
                Wortende $\rightarrow$ Endzustand entscheidet.
          \item \textbf{Historie}: jeder Schritt gespeichert
                $\rightarrow$ vor- und zurückspringen.
          \item Tabelle und Diagramm zeigen denselben Schritt, immer synchron.
        \end{itemize}
  }
\end{fullframe}
